\clearpage
\homeworktitle{Решения и комментарии}{Версия для учителя}

\begin{teacheranswer}[Задание 1]
  Из числителя получаем \(x=2\) и \(x=-3\). При \(x=2\) знаменатель
  равен нулю, поэтому это значение запрещено. Ответ:
  \[
    \boxed{x=-3}.
  \]
\end{teacheranswer}

\begin{teacheranswer}[Задание 2]
  \[
    2|x|-x=
    \begin{cases}
      x,&x\ge0,\\
      -3x,&x<0.
    \end{cases}
  \]
  При \(b<0\) корней нет, при \(b=0\) один корень, при \(b>0\) два.
\end{teacheranswer}

\section*{Решение задания 3}

\begin{teacheranswer}[Числитель]
  \[
    2|x|-x=a+1.
  \]
  При \(x\ge0\): \(x_1=a+1\), требуется \(a\ge-1\).
  При \(x<0\): \(x_2=-\dfrac{a+1}{3}\), требуется \(a>-1\).
  Два различных нуля возникают при \(a>-1\).
\end{teacheranswer}

\begin{teacheranswer}[Знаменатель и ответ]
  \[
    D(x_1)=(a+1)^2-(a+1)-a=a^2,
  \]
  поэтому исключается \(a=0\).
  \[
    D(x_2)=\frac{(a+1)^2}{9}+\frac{a+1}{3}-a
           =\frac{(a-2)^2}{9},
  \]
  поэтому исключается \(a=2\). Ответ:
  \[
    \boxed{a\in(-1,+\infty)\setminus\{0,2\}}.
  \]
\end{teacheranswer}

\section*{Решение задания 4}

\begin{teacheranswer}[Числитель]
  Из \(3|x|-x=a+2\) получаем
  \[
    x_1=\frac{a+2}{2},
    \qquad
    x_2=-\frac{a+2}{4}.
  \]
  Оба нуля существуют и различны при \(a>-2\).
\end{teacheranswer}

\begin{teacheranswer}[Знаменатель и ответ]
  \[
    D(x_1)=\frac{a(a-2)}4,
    \qquad
    D(x_2)=\frac{(a-2)(a-6)}{16}.
  \]
  Исключаются \(a=0,2,6\). Ответ:
  \[
    \boxed{a\in(-2,+\infty)\setminus\{0,2,6\}}.
  \]
\end{teacheranswer}

\begin{criterionbox}
  Значение \(a=2\) возникает при проверке обоих кандидатов, но в итоговом
  множестве исключается один раз.
\end{criterionbox}

\section*{Задания 5 и 6}

\begin{teacheranswer}[Задание 5]
  Условие \(a>-2\) правильно описывает только число нулей числителя. Решение
  становится неполным в момент, когда ученик не проверяет эти нули в
  знаменателе. Правильный ответ:
  \[
    \boxed{a\in(-2,+\infty)\setminus\{-1,0,3,8\}}.
  \]
\end{teacheranswer}

\begin{teacheranswer}[Задание 6]
  \[
    \begin{array}{c|c}
      a<-2 & 0\text{ корней},\\
      a=-2 & 1\text{ корень},\\
      a\in\{-1,0,3,8\} & 1\text{ корень},\\
      a>-2,\ a\notin\{-1,0,3,8\} & 2\text{ корня}.
    \end{array}
  \]
\end{teacheranswer}

\section*{Решение задания 7*}

\begin{teacheranswer}[Исследуем числитель]
  При \(x\ge a\):
  \[
    2(x-a)-x-2=0
    \quad\Longrightarrow\quad
    x_1=2a+2,
    \qquad a\ge-2.
  \]
  При \(x<a\):
  \[
    2(a-x)-x-2=0
    \quad\Longrightarrow\quad
    x_2=\frac{2a-2}{3},
    \qquad a>-2.
  \]
  Два различных нуля возникают при \(a>-2\).
\end{teacheranswer}

\begin{teacheranswer}[Проверяем знаменатель]
  Знаменатель равен нулю при \(x=\pm1\).
  Для \(x_1=2a+2\) получаем \(a=-\dfrac32,-\dfrac12\).
  Для \(x_2=\dfrac{2a-2}{3}\) получаем \(a=-\dfrac12,\dfrac52\).
  Ответ:
  \[
    \boxed{
      a\in(-2,+\infty)\setminus
      \left\{-\frac32,-\frac12,\frac52\right\}
    }.
  \]
\end{teacheranswer}

\begin{resultbox}[Краткие ответы]
  \begin{enumerate}
    \item \(x=-3\).
    \item \(b<0\): 0; \(b=0\): 1; \(b>0\): 2.
    \item \(a\in(-1,+\infty)\setminus\{0,2\}\).
    \item \(a\in(-2,+\infty)\setminus\{0,2,6\}\).
    \item \(a\in(-2,+\infty)\setminus\{-1,0,3,8\}\).
    \item \(0,1,1,2\) корня по строкам таблицы.
    \item \(a\in(-2,+\infty)\setminus
      \left\{-\dfrac32,-\dfrac12,\dfrac52\right\}\).
  \end{enumerate}
\end{resultbox}
